\section{Conclusion}

The Hyperledger Cacti project, with its ambitious goal of enabling interoperability across diverse DLTs, stands at a critical juncture where enhancing its accessibility and maintainability will significantly accelerate its adoption and community growth. This mentorship offers a unique opportunity to directly address key areas of improvement, transforming identified pain points into clear pathways for new contributors and streamlined development for maintainers.

My proposed phased plan, grounded in meticulous pre-application research and aligned with the project's strategic cleanup initiative, targets documentation fragmentation, developer onboarding friction, technical debt reduction, and CI/CD efficiency. Through the systematic execution of these phases, I am confident that I can deliver tangible, measurable improvements that not only benefit the Cacti codebase directly but also cultivate a more welcoming and productive environment for future contributions.

I am deeply enthusiastic about the prospect of contributing to Hyperledger Cacti and collaborating with its experienced maintainers and the wider LF Decentralized Trust community. My technical skills, commitment to open-source best practices, and proven ability to tackle complex challenges make me an ideal candidate to drive this initiative forward. I am eager to begin this journey and contribute to the continued success and growth of Hyperledger Cacti.
